\documentclass[12pt]{article}
\usepackage[utf8]{inputenc}
\usepackage[T1]{fontenc}
\usepackage{hyperref}
\title{Malecite-Passamaquoddy Language Resources for K-12 Curriculum Development}
\author{Compiled by Claude Sonnet 5.5 (claude-sonnet-5-5) for Andr\'as Kornai}
\date{}
\begin{document}
\maketitle
\section{Introduction}
Malecite-Passamaquoddy (Peskotomuhkati / Wolastoqey) is an Algic language of Canada and the United States. The strongest documented resource found is the 2008 dictionary of Francis and Leavitt, with an online version at the Peskotomuhkati-Wolastoqey Language Portal.
\section{Primary Educational Resources}
No dedicated textbooks or school programmes were found in this search.
\section{Linguistic and Reference Materials}
\begin{itemize}
\item A Passamaquoddy-Maliseet Dictionary (University of Maine Press, 2008, about 18,000 entries; Francis and Leavitt): \url{https://umaine.edu/umpress/?p=15}
\item UBC Relational Lexicography page for the dictionary: \url{https://knowledgebase.arts.ubc.ca/a-passamaquoddy-maliseet-dictionary-peskotomuhkati-wolastoqewi-latuwewakon}
\item Online version (UBC knowledgebase; portal with audio, example sentences and verb conjugations): \url{https://knowledgebase.arts.ubc.ca/?p=5407}
\end{itemize}
\section{Religious and Community Resources}
None were identified in these results.
\section{Digital and Audio Resources}
The online dictionary reportedly includes native-speaker audio for entries and links to video recordings; see the online version above.
\section{Curriculum Development Context}
The dictionary and portal give a strong lexical base, including audio. Pedagogical materials were not located and should be sought directly from the communities.
\end{document}